\section{Two Gnosticisms}

Evola remarks in his essay on Steiner\footnote{\url{https://gornahoor.net/library/KarmaAndReincarnation.pdf}} that Theosophy is actually a counterfeit doctrine. That is, the validity is true insofar as it goes, but the practical effect is actually to blunt the true dogma which it vaguely apprehends. Thus, modern ``reincarnation'' belief actually tends to inhibit the stark clarity which alone makes possible the choice between perdition and salvation, devolving (therefore) naturally to perdition. Here, as always, Evola speaks as a man who can discern without error the practical effects of abstraction. In another context, this might be known as a ``spiritual gift'':

\begin{quotex}
In fact, according to the Theosophical views, the ``gods'' and the adepts would be beings who had gone further ahead in ``evolution''; the animals, ``our younger brothers'', less ``advanced''. But it will be a question of time: everyone will reach the door, those who are further ahead ``sacrificing themselves'' for the others; and the varieties of karma will have served only as instrument to ``universal progress''. As is clear, all that can only be considered as a digressing and distorted addition of Theosophy to the authentic notion of karma. It should therefore not cause surprise if this notion often passes from the plane of a transcendental realism to a more or less Philistine moralism, becoming a type of sword of Damocles suspended over the head of whoever does not conform himself to the ``laws of evolution'' and to the related altruistic, humanistic, egalitarian, vegetarian, feminist, etc. corollaries professed by the movement. With that, even the practical value, the liberating potentiality of this teaching, which we already mentioned, must be lost completely.
\end{quotex}


This astounding claim hits very close to home for the ``modern'' person, whether New Age or Christian. Evola is essentially claiming that a counterfeit is actually worse than something (or anything) else, because an anti-Christ pretends or masquerades, blotting out and often discrediting the real article. Christ had a particular aversion to Pharisees, as to those who camp in front of the narrow door, with their backs to it, not intending to enter. The reason is simple --- it blocks the door.

Something similar happens in most versions of Christian universalism and modern Gnosticism and even in Western political history, processes which (like DNA) fray as they replicate until the original point is lost, and all hope of gaining truth or transcendence from the teaching has been vampirized.

Any religion, myth, story, government, family, narrative or group can serve the function of passing on a shabby truth whose point is forgotten --- this is no exclusive monopoly of the Catholic Church. It comes to pass that in the age of clay mixed with iron at the feet of King Nebuchadnezzar's statue, the atomistic individual is drawn face to face with inescapable questions about the true nature of the self, if he can see them. How can the Westerner see it?

To quote Rosenstock-Huessy\footnote{\url{http://www.argobooks.org/english/the_christian_future.html}} once more, we find him proclaiming the era of the individual:

\begin{quotex}
Today we are living through the agonies of transition to the third epoch. We have yet to establish Man, the great singular of humanity, in one household, over the plurality of races, classes, and and age groups. This will be the center of struggle in the future. They pose the questions the Third Millennium will have to answer... the State is on the defensive because it is inadequate for the needs of the coming age. The theme of future history will be not territorial nor political, but social.
\end{quotex}


Besides suggesting that we contract space and lengthen time (in opposition to the two modern tendencies to do the opposite), Huessy had few answers, only observations. Like most moderns, he had neglected Plato. Modernism is a ``hell'' of conflicting bodies.

Guenon and Evola and others have noted the ``relativity'' of the heavenly pattern as it relates to times and places, in which Heaven utilizes \emph{upaya} or skillful means to bring about the desired result to oppose chaos, lifting man out of his natural fate, even in the West\footnote{\url{http://www.gornahoor.net/?p=1224}}:

\begin{quotex}
But according to Anthony, only the rational man is in union with God (or \emph{Nous} in Platonic or Aristotelian terms), \emph{and reason is not at all common to all men}. Obviously, the modernist project of deriving human dignity from material and biological considerations has failed (it is unreasonable), with the resulting loss of cohesiveness in Western societies. So, if biologism cannot justify egalitarianism, then neither can Christianity, since full human dignity may only be virtual. Another way of putting it is to say that although in essence there may be a common divinity in men, it is not necessarily so in existence.
\end{quotex}


In a certain sense, materialism can be taken as a true condemnation of the futility of living purely empirically -- man is, as the atheist argued, ``a sack of protein and amino acid''. At its root, this is what Evola teaches -- without self-realization, the purely sensual human animal is condemned to undergo the ``hells''.

The teachings of Mouravieff are especially clear on this point, and his argument is that this was the original understanding of the Faith. As explained by George Heart (\textit{Dogmatic Faith \& Gnostic Vivifying Knowledge}):

\begin{quotex}
The belief that salvation can only occur after death in a nebulous paradise of some sort, is according to us a mortal error. Salvation is to be realized after a long and unceasing strife against sin -- properly understood -- on Earth and inside our physical bodies... it is (elsewhere) vaguely insinuated that a pious life on earth will save man's soul... (But) the mortal soul has to be integrated into the immortal soul -- not directly to God -- by means of a complex process of purification and sublimation. Salvation is the outcome of Man's unceasing, conscious and voluntary efforts... such is the meaning, necessity, and goal of the triple (tripartate) structure of man in Saint Paul. Let the dead (that is to say those in whom the mortal soul is still attached to the physical body) bury their dead.
\end{quotex}


You will also note here the agreement with Evola as to the possible danger of assimilating to something that is ``less than'' the true individual, which can occur in some varieties of mysticism. Mouraveiff further taught that reincarnation was ``true'' in the sense that man himself is an incomplete being -- he comes into existence when pre-existent divine sparks are united with psychic bodies that assume physical forms. Thus, the real ``man'' is not re-incarnating, only the imperishable spirit, which will attempt (again) to attain salvation. Salvation as popularly understood in an exoteric sense therefore occurs when the body, soul, and spirit are united under the spirit, which is called ``The Kingdom of God''. This complex process explains many of the vagaries, seeming contradictions, and deeper similiarities among the great traditions.

In the Christian tradition the \emph{pleroma} of Love is the falling away (maturing) of faith and hope as the fullness of Being is attained in this union inside of man, which puts man back as a judge over angels, and into the son-ship and daughter-hood of God. The uniqueness of Christ perfects paganism, but does not threaten it.

There is more than ``rapport'' here between Christianity in its purest form and traditionalism, there is almost total agreement. The modern Christian is generally ``concerned'' with the here \& now, and separates this from the ``afterlife'' \emph{in order to avoid this inconvenient duty of saving himself.} Here is Nicholas Gomez Davila\footnote{\url{http://don-colacho.blogspot.com/}} --

\begin{quotex}
Concerning himself intensely with his neighbor's condition allows the Christian to dissimulate to himself his doubts about the divinity of Christ and the existence of God. Charity can be the most subtle form of apostasy.
\end{quotex}


The Church should not be judged by its popular or populated forms, anymore than the New Age should be allowed to discredit Eastern tradition. The essential affinities between Evola's thought on salvation and reincarnation and Mouravieff's publication of what he argued was the esoteric ``pearl'' preserved within the Church, bear witness that Western tradition can be harmonized.


\flrightit{Posted on 2011-05-27 by Logres}